%%%%%%%%%%%%%%%%%%%%%%%%%%%%%%%%%%%%%%%%%%%%%%%%%%%%%%%%%%%%
% 8.11 NONPARAMETRIC STATISTICS
% The Composition of Advanced Mathematics
%%%%%%%%%%%%%%%%%%%%%%%%%%%%%%%%%%%%%%%%%%%%%%%%%%%%%%%%%%%%

\section{Nonparametric Statistics}
\label{sec:nonparametric-statistics}

\prerequisite{
Descriptive statistics, probability, random variables, sampling
distributions, confidence intervals, hypothesis testing, order
statistics, empirical distribution functions, categorical data, and
basic combinatorics.
}

\learningobjectives{
By the end of this section, the reader should be able to:
\begin{itemize}
    \item explain what is meant by nonparametric statistical inference;
    \item distinguish parametric and nonparametric procedures;
    \item distinguish distribution-free from assumption-free methods;
    \item explain the role of ranks and signs;
    \item construct and interpret empirical distribution functions;
    \item perform one-sample sign tests;
    \item construct confidence intervals for population medians using
          order statistics;
    \item perform Wilcoxon signed-rank tests;
    \item distinguish sign and signed-rank procedures;
    \item perform Mann--Whitney--Wilcoxon rank-sum tests;
    \item interpret the Mann--Whitney parameter carefully;
    \item perform Kruskal--Wallis tests;
    \item perform Friedman tests for blocked or repeated designs;
    \item compute Spearman rank correlation;
    \item compute Kendall rank correlation;
    \item understand tests based on Spearman and Kendall statistics;
    \item perform one-sample Kolmogorov--Smirnov tests conceptually;
    \item perform two-sample Kolmogorov--Smirnov tests conceptually;
    \item understand the Cramér--von Mises and Anderson--Darling ideas;
    \item understand permutation tests;
    \item distinguish permutation tests from rank tests;
    \item understand randomization tests in designed experiments;
    \item understand exact and asymptotic \(p\)-values;
    \item account for ties conceptually;
    \item understand rank-based confidence intervals;
    \item understand Hodges--Lehmann estimators conceptually;
    \item understand robustness and efficiency tradeoffs;
    \item recognize nonparametric density estimation;
    \item understand kernel density estimation conceptually;
    \item understand nonparametric regression conceptually;
    \item recognize bootstrap and resampling connections;
    \item choose appropriate nonparametric procedures for common study
          designs;
    \item connect nonparametric statistics with robust statistics,
          computational statistics, and statistical learning.
\end{itemize}
}

%%%%%%%%%%%%%%%%%%%%%%%%%%%%%%%%%%%%%%%%%%%%%%%%%%%%%%%%%%%%
\subsection{What Is Nonparametric Statistics?}
\label{subsec:what-is-nonparametric-statistics}
%%%%%%%%%%%%%%%%%%%%%%%%%%%%%%%%%%%%%%%%%%%%%%%%%%%%%%%%%%%%

Parametric statistical methods assume that data arise from a family of
probability distributions described by a finite-dimensional parameter
vector.

For example,

\[
\boxed{
X_i
\sim
N(\mu,\sigma^2)
}
\]

is a parametric model with parameters

\[
\mu
\]

and

\[
\sigma^2.
\]

Nonparametric statistics develops procedures that do not require such a
fully specified finite-dimensional distributional family.

%%%%%%%%%%%%%%%%%%%%%%%%%%%%%%%%%%%%%%%%%%%%%%%%%%%%%%%%%%%%
\subsection{Nonparametric Does Not Mean No Assumptions}
\label{subsec:nonparametric-not-assumption-free}
%%%%%%%%%%%%%%%%%%%%%%%%%%%%%%%%%%%%%%%%%%%%%%%%%%%%%%%%%%%%

Nonparametric methods are sometimes incorrectly described as
``assumption-free.''

They still require assumptions involving matters such as:

\[
\boxed{
\text{independence},
}
\]

\[
\boxed{
\text{exchangeability},
}
\]

\[
\boxed{
\text{continuity},
}
\]

\[
\boxed{
\text{symmetry},
}
\]

or

\[
\boxed{
\text{common distributional shape}.
}
\]

The exact assumptions depend on the procedure and the interpretation
desired.

%%%%%%%%%%%%%%%%%%%%%%%%%%%%%%%%%%%%%%%%%%%%%%%%%%%%%%%%%%%%
\subsection{Why Use Nonparametric Methods?}
\label{subsec:why-use-nonparametric}
%%%%%%%%%%%%%%%%%%%%%%%%%%%%%%%%%%%%%%%%%%%%%%%%%%%%%%%%%%%%

Nonparametric methods are useful when:

\[
\boxed{
\text{sample sizes are small},
}
\]

\[
\boxed{
\text{normality is questionable},
}
\]

\[
\boxed{
\text{outliers strongly affect parametric estimates},
}
\]

\[
\boxed{
\text{measurements are ordinal},
}
\]

or

\[
\boxed{
\text{the full population distribution is not specified}.
}
\]

They are also important because many statistical quantities themselves,
such as a distribution function, are infinite-dimensional objects.

%%%%%%%%%%%%%%%%%%%%%%%%%%%%%%%%%%%%%%%%%%%%%%%%%%%%%%%%%%%%
\subsection{Ranks}
\label{subsec:statistical-ranks}
%%%%%%%%%%%%%%%%%%%%%%%%%%%%%%%%%%%%%%%%%%%%%%%%%%%%%%%%%%%%

Suppose observations are

\[
X_1,\ldots,X_n.
\]

After sorting them,

\[
X_{(1)}
\leq
X_{(2)}
\leq
\cdots
\leq
X_{(n)},
\]

the smallest observation receives rank \(1\), the next rank \(2\), and
so forth.

Rank-based procedures replace the original measurement scale with
relative ordering information.

%%%%%%%%%%%%%%%%%%%%%%%%%%%%%%%%%%%%%%%%%%%%%%%%%%%%%%%%%%%%
\subsection{Ranks with Ties}
\label{subsec:ranks-with-ties}
%%%%%%%%%%%%%%%%%%%%%%%%%%%%%%%%%%%%%%%%%%%%%%%%%%%%%%%%%%%%

If several observations have the same value, a common convention assigns
each tied observation the average of the ranks they would have occupied.

For example, suppose values occupy ranks

\[
3,\ 4,\ 5.
\]

Each tied value receives rank

\[
\boxed{
\frac{3+4+5}{3}
=
4.
}
\]

Tie corrections may be required in exact variance formulas.

%%%%%%%%%%%%%%%%%%%%%%%%%%%%%%%%%%%%%%%%%%%%%%%%%%%%%%%%%%%%
\subsection{Signs}
\label{subsec:statistical-signs}
%%%%%%%%%%%%%%%%%%%%%%%%%%%%%%%%%%%%%%%%%%%%%%%%%%%%%%%%%%%%

Some procedures use only whether an observation lies above or below a
reference value.

For example, relative to hypothesized median \(m_0\),

\[
X_i-m_0
\]

has sign

\[
+
\]

if \(X_i>m_0\), and

\[
-
\]

if \(X_i<m_0\).

The magnitude of the deviation is ignored.

%%%%%%%%%%%%%%%%%%%%%%%%%%%%%%%%%%%%%%%%%%%%%%%%%%%%%%%%%%%%
\subsection{Empirical Distribution Function}
\label{subsec:edf-nonparametric}
%%%%%%%%%%%%%%%%%%%%%%%%%%%%%%%%%%%%%%%%%%%%%%%%%%%%%%%%%%%%

The empirical distribution function is

\[
\boxed{
F_n(x)
=
\frac1n
\sum_{i=1}^{n}
\mathbf{1}_{\{X_i\leq x\}}.
}
\]

It places probability mass

\[
\frac1n
\]

at each observed value.

The empirical distribution is one of the most fundamental nonparametric
estimators.

%%%%%%%%%%%%%%%%%%%%%%%%%%%%%%%%%%%%%%%%%%%%%%%%%%%%%%%%%%%%
\subsection{Properties of the Empirical Distribution Function}
\label{subsec:edf-properties}
%%%%%%%%%%%%%%%%%%%%%%%%%%%%%%%%%%%%%%%%%%%%%%%%%%%%%%%%%%%%

The EDF satisfies:

\[
\boxed{
0\leq F_n(x)\leq1,
}
\]

it is nondecreasing, and

\[
\boxed{
\lim_{x\to-\infty}F_n(x)=0,
}
\]

\[
\boxed{
\lim_{x\to\infty}F_n(x)=1.
}
\]

It is a right-continuous step function.

%%%%%%%%%%%%%%%%%%%%%%%%%%%%%%%%%%%%%%%%%%%%%%%%%%%%%%%%%%%%
\subsection{Glivenko--Cantelli Theorem Preview}
\label{subsec:glivenko-cantelli-preview}
%%%%%%%%%%%%%%%%%%%%%%%%%%%%%%%%%%%%%%%%%%%%%%%%%%%%%%%%%%%%

Under i.i.d. sampling from distribution function \(F\),

\[
\boxed{
\sup_x
|F_n(x)-F(x)|
\to0
}
\]

almost surely.

Thus the empirical distribution converges uniformly to the true
distribution.

This result provides a theoretical foundation for many nonparametric
methods.

%%%%%%%%%%%%%%%%%%%%%%%%%%%%%%%%%%%%%%%%%%%%%%%%%%%%%%%%%%%%
\subsection{One-Sample Sign Test}
\label{subsec:one-sample-sign-test}
%%%%%%%%%%%%%%%%%%%%%%%%%%%%%%%%%%%%%%%%%%%%%%%%%%%%%%%%%%%%

Suppose the goal is to test whether a continuous population has median

\[
m_0.
\]

Consider

\[
H_0:
m=m_0.
\]

Under the null, observations above and below \(m_0\) each have probability

\[
\frac12,
\]

assuming no probability mass occurs exactly at the median.

%%%%%%%%%%%%%%%%%%%%%%%%%%%%%%%%%%%%%%%%%%%%%%%%%%%%%%%%%%%%
\subsection{Sign-Test Statistic}
\label{subsec:sign-test-statistic}
%%%%%%%%%%%%%%%%%%%%%%%%%%%%%%%%%%%%%%%%%%%%%%%%%%%%%%%%%%%%

Let

\[
S
=
\#\{i:X_i>m_0\}.
\]

Ignoring observations exactly equal to \(m_0\), under the null,

\[
\boxed{
S
\sim
\operatorname{Binomial}
\left(
n,\frac12
\right).
}
\]

Thus the sign test reduces a location problem to a binomial test.

%%%%%%%%%%%%%%%%%%%%%%%%%%%%%%%%%%%%%%%%%%%%%%%%%%%%%%%%%%%%
\subsection{Worked Example: Sign Test}
\label{subsec:worked-sign-test}
%%%%%%%%%%%%%%%%%%%%%%%%%%%%%%%%%%%%%%%%%%%%%%%%%%%%%%%%%%%%

\begin{workedexamplebox}{Testing a Median}

Suppose \(12\) non-tied observations are compared with hypothesized
median

\[
m_0=50.
\]

Suppose

\[
9
\]

observations exceed \(50\), while

\[
3
\]

fall below \(50\).

Under

\[
H_0:m=50,
\]

\[
S
\sim
\operatorname{Binomial}
\left(
12,\frac12
\right).
\]

For a right-tailed alternative,

\[
\begin{answerbox}
\[
\boxed{
p\text{-value}
=
P(S\geq9)
=
\sum_{j=9}^{12}
\binom{12}{j}
\left(
\frac12
\right)^{12}.
}
\]
\end{answerbox}

\end{workedexamplebox}

%%%%%%%%%%%%%%%%%%%%%%%%%%%%%%%%%%%%%%%%%%%%%%%%%%%%%%%%%%%%
\subsection{Handling Ties in the Sign Test}
\label{subsec:sign-test-ties}
%%%%%%%%%%%%%%%%%%%%%%%%%%%%%%%%%%%%%%%%%%%%%%%%%%%%%%%%%%%%

Observations satisfying

\[
X_i=m_0
\]

contain no directional information.

A common exact procedure removes these ties and performs the binomial test
using only the remaining observations.

Thus the effective sample size may be smaller than the original sample
size.

%%%%%%%%%%%%%%%%%%%%%%%%%%%%%%%%%%%%%%%%%%%%%%%%%%%%%%%%%%%%
\subsection{Advantages of the Sign Test}
\label{subsec:sign-test-advantages}
%%%%%%%%%%%%%%%%%%%%%%%%%%%%%%%%%%%%%%%%%%%%%%%%%%%%%%%%%%%%

The sign test requires little distributional structure.

It is resistant to extreme observations because it uses only signs.

For example,

\[
51
\]

and

\[
5000
\]

both contribute one positive sign when testing against

\[
50.
\]

%%%%%%%%%%%%%%%%%%%%%%%%%%%%%%%%%%%%%%%%%%%%%%%%%%%%%%%%%%%%
\subsection{Limitation of the Sign Test}
\label{subsec:sign-test-limitation}
%%%%%%%%%%%%%%%%%%%%%%%%%%%%%%%%%%%%%%%%%%%%%%%%%%%%%%%%%%%%

Because the sign test ignores magnitudes, it can discard useful
information.

Thus it may have lower power than procedures that use both direction and
magnitude when stronger assumptions are valid.

%%%%%%%%%%%%%%%%%%%%%%%%%%%%%%%%%%%%%%%%%%%%%%%%%%%%%%%%%%%%
\subsection{Distribution-Free Confidence Interval for a Median}
\label{subsec:distribution-free-median-ci}
%%%%%%%%%%%%%%%%%%%%%%%%%%%%%%%%%%%%%%%%%%%%%%%%%%%%%%%%%%%%

Suppose

\[
X_{(1)}
\leq
\cdots
\leq
X_{(n)}
\]

are order statistics from a continuous distribution with median \(m\).

Since each observation lies below \(m\) with probability \(1/2\), binomial
probabilities can be used to select indices \(r\) and \(s\) such that

\[
\boxed{
[X_{(r)},X_{(s)}]
}
\]

has at least a desired confidence level for the population median.

This interval requires no assumption of normality.

%%%%%%%%%%%%%%%%%%%%%%%%%%%%%%%%%%%%%%%%%%%%%%%%%%%%%%%%%%%%
\subsection{Order-Statistic Logic for Median Intervals}
\label{subsec:median-order-statistic-logic}
%%%%%%%%%%%%%%%%%%%%%%%%%%%%%%%%%%%%%%%%%%%%%%%%%%%%%%%%%%%%

The event

\[
X_{(r)}
\leq
m
\leq
X_{(s)}
\]

is determined by the number of sample observations below the population
median.

Because that count has a binomial distribution with success probability

\[
\frac12,
\]

exact finite-sample coverage can be calculated.

%%%%%%%%%%%%%%%%%%%%%%%%%%%%%%%%%%%%%%%%%%%%%%%%%%%%%%%%%%%%
\subsection{Paired Nonparametric Procedures}
\label{subsec:paired-nonparametric}
%%%%%%%%%%%%%%%%%%%%%%%%%%%%%%%%%%%%%%%%%%%%%%%%%%%%%%%%%%%%

Suppose paired observations are

\[
(X_i,Y_i).
\]

Define differences

\[
\boxed{
D_i=X_i-Y_i.
}
\]

Nonparametric paired procedures then analyze the distribution of

\[
D_1,\ldots,D_n.
\]

The sign test uses only

\[
\operatorname{sign}(D_i).
\]

The Wilcoxon signed-rank test also uses the magnitude ranking of
\(|D_i|\).

%%%%%%%%%%%%%%%%%%%%%%%%%%%%%%%%%%%%%%%%%%%%%%%%%%%%%%%%%%%%
\subsection{Wilcoxon Signed-Rank Test}
\label{subsec:wilcoxon-signed-rank}
%%%%%%%%%%%%%%%%%%%%%%%%%%%%%%%%%%%%%%%%%%%%%%%%%%%%%%%%%%%%

The Wilcoxon signed-rank test is a rank-based procedure for paired or
one-sample location problems.

For differences

\[
D_1,\ldots,D_n,
\]

the procedure is:

\begin{enumerate}
    \item remove zero differences;
    \item compute \(|D_i|\);
    \item rank the absolute differences;
    \item restore the original signs;
    \item sum positive and negative signed ranks.
\end{enumerate}

%%%%%%%%%%%%%%%%%%%%%%%%%%%%%%%%%%%%%%%%%%%%%%%%%%%%%%%%%%%%
\subsection{Signed-Rank Statistics}
\label{subsec:signed-rank-statistics}
%%%%%%%%%%%%%%%%%%%%%%%%%%%%%%%%%%%%%%%%%%%%%%%%%%%%%%%%%%%%

Let

\[
R_i
\]

be the rank of

\[
|D_i|.
\]

Define

\[
\boxed{
W^+
=
\sum_{D_i>0}R_i,
}
\]

and

\[
\boxed{
W^-
=
\sum_{D_i<0}R_i.
}
\]

Since the ranks are

\[
1,\ldots,n
\]

when there are no ties,

\[
\boxed{
W^++W^-
=
\frac{
n(n+1)
}{2}.
}
\]

%%%%%%%%%%%%%%%%%%%%%%%%%%%%%%%%%%%%%%%%%%%%%%%%%%%%%%%%%%%%
\subsection{Worked Example: Signed Ranks}
\label{subsec:worked-signed-ranks}
%%%%%%%%%%%%%%%%%%%%%%%%%%%%%%%%%%%%%%%%%%%%%%%%%%%%%%%%%%%%

\begin{workedexamplebox}{Computing \(W^+\)}

Suppose the paired differences are

\[
-2,\ 1,\ 4,\ -3.
\]

Their absolute values are

\[
2,\ 1,\ 4,\ 3.
\]

The absolute ranks are

\[
2,\ 1,\ 4,\ 3.
\]

Restoring signs gives

\[
-2,\ +1,\ +4,\ -3.
\]

Therefore,

\[
W^+
=
1+4
=
5.
\]

Also,

\[
W^-
=
2+3
=
5.
\]

\begin{answerbox}
\[
\boxed{
W^+=5,
\qquad
W^-=5.
}
\]
\end{answerbox}

\end{workedexamplebox}

%%%%%%%%%%%%%%%%%%%%%%%%%%%%%%%%%%%%%%%%%%%%%%%%%%%%%%%%%%%%
\subsection{Signed-Rank Assumptions}
\label{subsec:signed-rank-assumptions}
%%%%%%%%%%%%%%%%%%%%%%%%%%%%%%%%%%%%%%%%%%%%%%%%%%%%%%%%%%%%

For the classical location interpretation, the differences are commonly
assumed to arise independently from a continuous distribution symmetric
about a location parameter.

Thus the signed-rank test requires more structure than the sign test.

It does not require normality.

%%%%%%%%%%%%%%%%%%%%%%%%%%%%%%%%%%%%%%%%%%%%%%%%%%%%%%%%%%%%
\subsection{Sign Test Versus Signed-Rank Test}
\label{subsec:sign-versus-signed-rank}
%%%%%%%%%%%%%%%%%%%%%%%%%%%%%%%%%%%%%%%%%%%%%%%%%%%%%%%%%%%%

The sign test uses only

\[
\boxed{
\text{direction}.
}
\]

The signed-rank test uses

\[
\boxed{
\text{direction}
+
\text{relative magnitude}.
}
\]

Therefore the signed-rank test can be more efficient when its symmetry
assumption is appropriate.

The sign test is more robust to asymmetric distributions.

%%%%%%%%%%%%%%%%%%%%%%%%%%%%%%%%%%%%%%%%%%%%%%%%%%%%%%%%%%%%
\subsection{Expected Signed-Rank Sum}
\label{subsec:expected-signed-rank}
%%%%%%%%%%%%%%%%%%%%%%%%%%%%%%%%%%%%%%%%%%%%%%%%%%%%%%%%%%%%

Under a symmetric null with no ties, positive and negative signs are
equally likely.

Since the total rank sum is

\[
\frac{
n(n+1)
}{2},
\]

the expected positive-rank sum is

\[
\boxed{
\mathbb{E}[W^+]
=
\frac{
n(n+1)
}{4}.
}
\]

%%%%%%%%%%%%%%%%%%%%%%%%%%%%%%%%%%%%%%%%%%%%%%%%%%%%%%%%%%%%
\subsection{Variance of the Signed-Rank Statistic}
\label{subsec:variance-signed-rank}
%%%%%%%%%%%%%%%%%%%%%%%%%%%%%%%%%%%%%%%%%%%%%%%%%%%%%%%%%%%%

Without ties,

\[
\boxed{
\operatorname{Var}(W^+)
=
\frac{
n(n+1)(2n+1)
}{
24
}.
}
\]

This leads to a normal approximation for sufficiently large samples.

Tie corrections alter the variance.

%%%%%%%%%%%%%%%%%%%%%%%%%%%%%%%%%%%%%%%%%%%%%%%%%%%%%%%%%%%%
\subsection{Normal Approximation to Signed-Rank}
\label{subsec:normal-approx-signed-rank}
%%%%%%%%%%%%%%%%%%%%%%%%%%%%%%%%%%%%%%%%%%%%%%%%%%%%%%%%%%%%

For sufficiently large \(n\),

\[
\boxed{
Z
=
\frac{
W^+
-
n(n+1)/4
}{
\sqrt{
n(n+1)(2n+1)/24
}
}
}
\]

is approximately standard normal under the null when no tie corrections
are needed.

A continuity correction may also be applied.

%%%%%%%%%%%%%%%%%%%%%%%%%%%%%%%%%%%%%%%%%%%%%%%%%%%%%%%%%%%%
\subsection{Two Independent Samples}
\label{subsec:two-independent-nonparametric}
%%%%%%%%%%%%%%%%%%%%%%%%%%%%%%%%%%%%%%%%%%%%%%%%%%%%%%%%%%%%

Suppose two independent samples are

\[
X_1,\ldots,X_{n_1}
\]

and

\[
Y_1,\ldots,Y_{n_2}.
\]

A common nonparametric question is whether the two populations have the
same distribution.

The Mann--Whitney--Wilcoxon procedure is a rank-based method for this
setting.

%%%%%%%%%%%%%%%%%%%%%%%%%%%%%%%%%%%%%%%%%%%%%%%%%%%%%%%%%%%%
\subsection{Wilcoxon Rank-Sum Test}
\label{subsec:wilcoxon-rank-sum}
%%%%%%%%%%%%%%%%%%%%%%%%%%%%%%%%%%%%%%%%%%%%%%%%%%%%%%%%%%%%

Combine all

\[
N=n_1+n_2
\]

observations and rank them together.

Let

\[
R_1
\]

be the sum of ranks belonging to sample \(1\).

Under the null hypothesis of identical continuous distributions, the
allocation of ranks between groups is exchangeable.

%%%%%%%%%%%%%%%%%%%%%%%%%%%%%%%%%%%%%%%%%%%%%%%%%%%%%%%%%%%%
\subsection{Expected Rank Sum}
\label{subsec:expected-rank-sum}
%%%%%%%%%%%%%%%%%%%%%%%%%%%%%%%%%%%%%%%%%%%%%%%%%%%%%%%%%%%%

Under the null and without ties,

\[
\boxed{
\mathbb{E}[R_1]
=
\frac{
n_1(N+1)
}{2}.
}
\]

The variance is

\[
\boxed{
\operatorname{Var}(R_1)
=
\frac{
n_1n_2(N+1)
}{
12
}.
}
\]

This leads to a large-sample normal approximation.

%%%%%%%%%%%%%%%%%%%%%%%%%%%%%%%%%%%%%%%%%%%%%%%%%%%%%%%%%%%%
\subsection{Mann--Whitney \(U\) Statistic}
\label{subsec:mann-whitney-u}
%%%%%%%%%%%%%%%%%%%%%%%%%%%%%%%%%%%%%%%%%%%%%%%%%%%%%%%%%%%%

The rank-sum statistic can be transformed into the Mann--Whitney
statistic:

\[
\boxed{
U_1
=
R_1
-
\frac{
n_1(n_1+1)
}{2}.
}
\]

Similarly,

\[
\boxed{
U_2
=
n_1n_2-U_1.
}
\]

Thus the Wilcoxon rank-sum and Mann--Whitney tests are equivalent
formulations.

%%%%%%%%%%%%%%%%%%%%%%%%%%%%%%%%%%%%%%%%%%%%%%%%%%%%%%%%%%%%
\subsection{Pairwise Interpretation of \(U\)}
\label{subsec:mann-whitney-pairwise}
%%%%%%%%%%%%%%%%%%%%%%%%%%%%%%%%%%%%%%%%%%%%%%%%%%%%%%%%%%%%

The statistic \(U\) can be interpreted through pairwise comparisons.

Ignoring ties,

\[
U_1
\]

counts the number of pairs

\[
(X_i,Y_j)
\]

for which

\[
X_i>Y_j.
\]

With ties, half-credit is commonly assigned.

Thus a normalized version estimates

\[
\boxed{
P(X>Y)
+
\frac12P(X=Y).
}
\]

%%%%%%%%%%%%%%%%%%%%%%%%%%%%%%%%%%%%%%%%%%%%%%%%%%%%%%%%%%%%
\subsection{Mann--Whitney Does Not Automatically Test Medians}
\label{subsec:mann-whitney-not-median}
%%%%%%%%%%%%%%%%%%%%%%%%%%%%%%%%%%%%%%%%%%%%%%%%%%%%%%%%%%%%

The Mann--Whitney test is often described as a test of medians.

That interpretation requires additional assumptions, such as similarly
shaped distributions differing primarily by location.

More generally, the procedure detects distributional differences related
to relative ordering.

Thus:

\[
\boxed{
\text{Mann--Whitney}
\neq
\text{universal median test}.
}
\]

%%%%%%%%%%%%%%%%%%%%%%%%%%%%%%%%%%%%%%%%%%%%%%%%%%%%%%%%%%%%
\subsection{Worked Example: Rank Sum}
\label{subsec:worked-rank-sum}
%%%%%%%%%%%%%%%%%%%%%%%%%%%%%%%%%%%%%%%%%%%%%%%%%%%%%%%%%%%%

\begin{workedexamplebox}{Two Independent Groups}

Suppose group \(A\) contains

\[
2,\ 5,\ 8
\]

and group \(B\) contains

\[
1,\ 4,\ 6.
\]

The combined ordered observations are

\[
1_B,\ 2_A,\ 4_B,\ 5_A,\ 6_B,\ 8_A.
\]

Thus group \(A\) receives ranks

\[
2,\ 4,\ 6.
\]

Therefore,

\[
R_A
=
2+4+6
=
12.
\]

The corresponding \(U\)-statistic is

\[
U_A
=
12
-
\frac{
3(4)
}{2}
=
6.
\]

\begin{answerbox}
\[
\boxed{
R_A=12,
\qquad
U_A=6.
}
\]
\end{answerbox}

\end{workedexamplebox}

%%%%%%%%%%%%%%%%%%%%%%%%%%%%%%%%%%%%%%%%%%%%%%%%%%%%%%%%%%%%
\subsection{Normal Approximation for the Rank-Sum Test}
\label{subsec:rank-sum-normal-approx}
%%%%%%%%%%%%%%%%%%%%%%%%%%%%%%%%%%%%%%%%%%%%%%%%%%%%%%%%%%%%

For sufficiently large samples and no ties,

\[
\boxed{
Z
=
\frac{
R_1
-
n_1(N+1)/2
}{
\sqrt{
n_1n_2(N+1)/12
}
}
}
\]

is approximately standard normal under the null.

Tie corrections and continuity corrections may be incorporated when
appropriate.

%%%%%%%%%%%%%%%%%%%%%%%%%%%%%%%%%%%%%%%%%%%%%%%%%%%%%%%%%%%%
\subsection{Kruskal--Wallis Test}
\label{subsec:kruskal-wallis-test}
%%%%%%%%%%%%%%%%%%%%%%%%%%%%%%%%%%%%%%%%%%%%%%%%%%%%%%%%%%%%

The Kruskal--Wallis procedure extends the rank-sum idea to

\[
k
\]

independent groups.

All observations are ranked jointly.

Let

\[
R_i
\]

denote the sum of ranks in group \(i\), and

\[
n_i
\]

its sample size.

%%%%%%%%%%%%%%%%%%%%%%%%%%%%%%%%%%%%%%%%%%%%%%%%%%%%%%%%%%%%
\subsection{Kruskal--Wallis Statistic}
\label{subsec:kruskal-wallis-statistic}
%%%%%%%%%%%%%%%%%%%%%%%%%%%%%%%%%%%%%%%%%%%%%%%%%%%%%%%%%%%%

Without ties, the statistic is

\[
\boxed{
H
=
\frac{
12
}{
N(N+1)
}
\sum_{i=1}^{k}
\frac{
R_i^2
}{
n_i
}
-
3(N+1).
}
\]

Under the null and suitable large-sample conditions,

\[
\boxed{
H
\approx
\chi^2_{k-1}.
}
\]

Tie corrections are used when ranks are tied.

%%%%%%%%%%%%%%%%%%%%%%%%%%%%%%%%%%%%%%%%%%%%%%%%%%%%%%%%%%%%
\subsection{Kruskal--Wallis Hypothesis}
\label{subsec:kruskal-wallis-hypothesis}
%%%%%%%%%%%%%%%%%%%%%%%%%%%%%%%%%%%%%%%%%%%%%%%%%%%%%%%%%%%%

A broad null hypothesis is that the \(k\) independent samples arise from
the same distribution.

Under additional common-shape assumptions, the procedure can be
interpreted as testing equality of location parameters.

It should not automatically be described as a test of equal medians.

%%%%%%%%%%%%%%%%%%%%%%%%%%%%%%%%%%%%%%%%%%%%%%%%%%%%%%%%%%%%
\subsection{Worked Example: Kruskal--Wallis Structure}
\label{subsec:worked-kruskal-wallis}
%%%%%%%%%%%%%%%%%%%%%%%%%%%%%%%%%%%%%%%%%%%%%%%%%%%%%%%%%%%%

\begin{workedexamplebox}{Three Groups}

Suppose three independent groups have sizes

\[
n_1=4,
\qquad
n_2=4,
\qquad
n_3=4.
\]

After ranking all \(12\) observations, suppose the rank sums are

\[
R_1=18,
\qquad
R_2=26,
\qquad
R_3=34.
\]

Then

\[
H
=
\frac{
12
}{
12(13)
}
\left[
\frac{18^2}{4}
+
\frac{26^2}{4}
+
\frac{34^2}{4}
\right]
-
3(13).
\]

\begin{answerbox}
\[
\boxed{
H
=
\frac1{13}
\left[
81+169+289
\right]
-
39.
}
\]
\end{answerbox}

The resulting statistic is compared with a chi-square distribution with

\[
\boxed{
2
}
\]

degrees of freedom.

\end{workedexamplebox}

%%%%%%%%%%%%%%%%%%%%%%%%%%%%%%%%%%%%%%%%%%%%%%%%%%%%%%%%%%%%
\subsection{Post Hoc Comparisons After Kruskal--Wallis}
\label{subsec:posthoc-kruskal-wallis}
%%%%%%%%%%%%%%%%%%%%%%%%%%%%%%%%%%%%%%%%%%%%%%%%%%%%%%%%%%%%

A significant Kruskal--Wallis test indicates that not all group
distributions are alike.

It does not identify the differing groups.

Possible follow-up procedures compare average ranks or pairwise rank-sum
statistics with multiplicity adjustment.

Common approaches include:

\[
\boxed{
\text{Dunn-type comparisons},
}
\]

or pairwise rank-sum tests with

\[
\boxed{
\text{Bonferroni}
}
\]

or other corrections.

%%%%%%%%%%%%%%%%%%%%%%%%%%%%%%%%%%%%%%%%%%%%%%%%%%%%%%%%%%%%
\subsection{Friedman Test}
\label{subsec:friedman-test}
%%%%%%%%%%%%%%%%%%%%%%%%%%%%%%%%%%%%%%%%%%%%%%%%%%%%%%%%%%%%

The Friedman test is a rank-based method for blocked or repeated-measures
data.

Suppose

\[
b
\]

blocks each receive

\[
k
\]

treatments.

Within each block, rank the \(k\) treatment responses.

Let

\[
R_j
\]

be the sum of ranks for treatment \(j\).

%%%%%%%%%%%%%%%%%%%%%%%%%%%%%%%%%%%%%%%%%%%%%%%%%%%%%%%%%%%%
\subsection{Friedman Statistic}
\label{subsec:friedman-statistic}
%%%%%%%%%%%%%%%%%%%%%%%%%%%%%%%%%%%%%%%%%%%%%%%%%%%%%%%%%%%%

Without ties, the Friedman statistic is

\[
\boxed{
Q
=
\frac{
12
}{
bk(k+1)
}
\sum_{j=1}^{k}
R_j^2
-
3b(k+1).
}
\]

Under the null and suitable sample sizes,

\[
\boxed{
Q
\approx
\chi^2_{k-1}.
}
\]

%%%%%%%%%%%%%%%%%%%%%%%%%%%%%%%%%%%%%%%%%%%%%%%%%%%%%%%%%%%%
\subsection{Friedman Versus Kruskal--Wallis}
\label{subsec:friedman-versus-kruskal}
%%%%%%%%%%%%%%%%%%%%%%%%%%%%%%%%%%%%%%%%%%%%%%%%%%%%%%%%%%%%

Kruskal--Wallis is designed for

\[
\boxed{
\text{independent groups}.
}
\]

Friedman is designed for

\[
\boxed{
\text{blocked or repeated observations}.
}
\]

Thus the distinction parallels:

\[
\boxed{
\text{one-way independent ANOVA}
}
\]

versus

\[
\boxed{
\text{randomized-block or repeated-measures analysis}.
}
\]

%%%%%%%%%%%%%%%%%%%%%%%%%%%%%%%%%%%%%%%%%%%%%%%%%%%%%%%%%%%%
\subsection{Worked Example: Friedman Ranking}
\label{subsec:worked-friedman-ranking}
%%%%%%%%%%%%%%%%%%%%%%%%%%%%%%%%%%%%%%%%%%%%%%%%%%%%%%%%%%%%

\begin{workedexamplebox}{Three Treatments Within a Block}

Suppose one block produces treatment responses

\[
8,\ 5,\ 11.
\]

Within that block, the ranks are

\[
2,\ 1,\ 3.
\]

These within-block ranks are accumulated across blocks.

\begin{answerbox}
\[
\boxed{
\text{Ranks within each block, not across all observations globally.}
}
\]
\end{answerbox}

\end{workedexamplebox}

%%%%%%%%%%%%%%%%%%%%%%%%%%%%%%%%%%%%%%%%%%%%%%%%%%%%%%%%%%%%
\subsection{Rank Correlation}
\label{subsec:rank-correlation}
%%%%%%%%%%%%%%%%%%%%%%%%%%%%%%%%%%%%%%%%%%%%%%%%%%%%%%%%%%%%

Pearson correlation measures linear association between numerical
variables.

Rank correlations measure monotonic association using relative ordering.

The two most common rank correlations are:

\[
\boxed{
\text{Spearman's }\rho_s
}
\]

and

\[
\boxed{
\text{Kendall's }\tau.
}
\]

%%%%%%%%%%%%%%%%%%%%%%%%%%%%%%%%%%%%%%%%%%%%%%%%%%%%%%%%%%%%
\subsection{Spearman Rank Correlation}
\label{subsec:spearman-rank-correlation}
%%%%%%%%%%%%%%%%%%%%%%%%%%%%%%%%%%%%%%%%%%%%%%%%%%%%%%%%%%%%

Spearman's rank correlation is the Pearson correlation between the rank
variables.

If

\[
R_i
\]

is the rank of \(X_i\) and

\[
S_i
\]

the rank of \(Y_i\), then

\[
\boxed{
\rho_s
=
\operatorname{Corr}(R,S).
}
\]

The sample version is calculated from observed ranks.

%%%%%%%%%%%%%%%%%%%%%%%%%%%%%%%%%%%%%%%%%%%%%%%%%%%%%%%%%%%%
\subsection{Spearman Formula Without Ties}
\label{subsec:spearman-no-ties}
%%%%%%%%%%%%%%%%%%%%%%%%%%%%%%%%%%%%%%%%%%%%%%%%%%%%%%%%%%%%

When there are no ties, define

\[
d_i
=
R_i-S_i.
\]

Then

\[
\boxed{
r_s
=
1
-
\frac{
6\sum_{i=1}^{n}d_i^2
}{
n(n^2-1)
}.
}
\]

This shortcut is equivalent to Pearson correlation of ranks when ranks
are distinct.

%%%%%%%%%%%%%%%%%%%%%%%%%%%%%%%%%%%%%%%%%%%%%%%%%%%%%%%%%%%%
\subsection{Worked Example: Spearman Correlation}
\label{subsec:worked-spearman}
%%%%%%%%%%%%%%%%%%%%%%%%%%%%%%%%%%%%%%%%%%%%%%%%%%%%%%%%%%%%

\begin{workedexamplebox}{Rank Agreement}

Suppose the ranks are

\[
R=(1,2,3,4)
\]

and

\[
S=(1,3,2,4).
\]

Then

\[
d=(0,-1,1,0).
\]

Therefore,

\[
\sum d_i^2
=
2.
\]

Thus,

\[
r_s
=
1
-
\frac{
6(2)
}{
4(4^2-1)
}.
\]

\begin{answerbox}
\[
\boxed{
r_s
=
1-\frac{12}{60}
=
0.8.
}
\]
\end{answerbox}

\end{workedexamplebox}

%%%%%%%%%%%%%%%%%%%%%%%%%%%%%%%%%%%%%%%%%%%%%%%%%%%%%%%%%%%%
\subsection{Interpretation of Spearman Correlation}
\label{subsec:spearman-interpretation}
%%%%%%%%%%%%%%%%%%%%%%%%%%%%%%%%%%%%%%%%%%%%%%%%%%%%%%%%%%%%

Spearman correlation measures monotonic association.

A relationship may be strongly monotone but nonlinear.

For example,

\[
Y=e^X
\]

can produce a rank correlation near \(1\) even though the relationship
is not linear.

Thus Spearman correlation is less tied to linear functional form than
Pearson correlation.

%%%%%%%%%%%%%%%%%%%%%%%%%%%%%%%%%%%%%%%%%%%%%%%%%%%%%%%%%%%%
\subsection{Kendall's Tau}
\label{subsec:kendalls-tau}
%%%%%%%%%%%%%%%%%%%%%%%%%%%%%%%%%%%%%%%%%%%%%%%%%%%%%%%%%%%%

Kendall's tau is based on concordant and discordant observation pairs.

Consider two observations

\[
(X_i,Y_i)
\]

and

\[
(X_j,Y_j).
\]

They are concordant if the order of the \(X\)-values agrees with the
order of the \(Y\)-values.

They are discordant if the orders disagree.

%%%%%%%%%%%%%%%%%%%%%%%%%%%%%%%%%%%%%%%%%%%%%%%%%%%%%%%%%%%%
\subsection{Kendall Tau Without Ties}
\label{subsec:kendall-no-ties}
%%%%%%%%%%%%%%%%%%%%%%%%%%%%%%%%%%%%%%%%%%%%%%%%%%%%%%%%%%%%

Let

\[
C
\]

be the number of concordant pairs and

\[
D
\]

the number of discordant pairs.

There are

\[
\binom{n}{2}
\]

pairs.

Without ties,

\[
\boxed{
\tau
=
\frac{
C-D
}{
\binom{n}{2}
}.
}
\]

Thus,

\[
-1\leq\tau\leq1.
\]

%%%%%%%%%%%%%%%%%%%%%%%%%%%%%%%%%%%%%%%%%%%%%%%%%%%%%%%%%%%%
\subsection{Worked Example: Kendall Tau}
\label{subsec:worked-kendall}
%%%%%%%%%%%%%%%%%%%%%%%%%%%%%%%%%%%%%%%%%%%%%%%%%%%%%%%%%%%%

\begin{workedexamplebox}{Concordant and Discordant Pairs}

Suppose \(n=4\), so there are

\[
\binom42=6
\]

pairs.

Suppose

\[
C=5
\]

and

\[
D=1.
\]

Then

\[
\tau
=
\frac{5-1}{6}.
\]

\begin{answerbox}
\[
\boxed{
\tau
=
\frac23.
}
\]
\end{answerbox}

\end{workedexamplebox}

%%%%%%%%%%%%%%%%%%%%%%%%%%%%%%%%%%%%%%%%%%%%%%%%%%%%%%%%%%%%
\subsection{Kendall Tau with Ties}
\label{subsec:kendall-ties}
%%%%%%%%%%%%%%%%%%%%%%%%%%%%%%%%%%%%%%%%%%%%%%%%%%%%%%%%%%%%

Several versions of Kendall's statistic exist.

A commonly used adjustment is

\[
\tau_b,
\]

which accounts for ties in both variables.

Therefore software output should be interpreted according to the exact
variant reported.

%%%%%%%%%%%%%%%%%%%%%%%%%%%%%%%%%%%%%%%%%%%%%%%%%%%%%%%%%%%%
\subsection{Spearman Versus Kendall}
\label{subsec:spearman-versus-kendall}
%%%%%%%%%%%%%%%%%%%%%%%%%%%%%%%%%%%%%%%%%%%%%%%%%%%%%%%%%%%%

Both measure monotonic association.

Spearman correlation uses rank values directly.

Kendall correlation uses pairwise concordance.

Kendall's tau often has a direct probability interpretation through the
difference between concordance and discordance probabilities.

%%%%%%%%%%%%%%%%%%%%%%%%%%%%%%%%%%%%%%%%%%%%%%%%%%%%%%%%%%%%
\subsection{Testing Rank Correlations}
\label{subsec:testing-rank-correlations}
%%%%%%%%%%%%%%%%%%%%%%%%%%%%%%%%%%%%%%%%%%%%%%%%%%%%%%%%%%%%

The null hypothesis may be

\[
\boxed{
H_0:
\text{no rank association}.
}
\]

For small samples without ties, exact permutation distributions can be
used.

For larger samples, normal or other asymptotic approximations are
available.

The precise null interpretation differs between Spearman and Kendall
statistics.

%%%%%%%%%%%%%%%%%%%%%%%%%%%%%%%%%%%%%%%%%%%%%%%%%%%%%%%%%%%%
\subsection{One-Sample Kolmogorov--Smirnov Test}
\label{subsec:one-sample-ks}
%%%%%%%%%%%%%%%%%%%%%%%%%%%%%%%%%%%%%%%%%%%%%%%%%%%%%%%%%%%%

Suppose the null hypothesis specifies a continuous distribution

\[
F_0.
\]

The one-sample Kolmogorov--Smirnov statistic is

\[
\boxed{
D_n
=
\sup_x
|F_n(x)-F_0(x)|.
}
\]

It measures the largest vertical distance between the empirical CDF and
the hypothesized CDF.

%%%%%%%%%%%%%%%%%%%%%%%%%%%%%%%%%%%%%%%%%%%%%%%%%%%%%%%%%%%%
\subsection{Interpretation of the KS Statistic}
\label{subsec:ks-interpretation}
%%%%%%%%%%%%%%%%%%%%%%%%%%%%%%%%%%%%%%%%%%%%%%%%%%%%%%%%%%%%

If

\[
D_n
\]

is small, the empirical distribution lies close to \(F_0\).

If

\[
D_n
\]

is large, some part of the empirical distribution deviates substantially
from the hypothesized model.

Thus the KS test compares entire distributions rather than only one
moment such as the mean.

%%%%%%%%%%%%%%%%%%%%%%%%%%%%%%%%%%%%%%%%%%%%%%%%%%%%%%%%%%%%
\subsection{Important KS Qualification}
\label{subsec:ks-parameter-estimation-qualification}
%%%%%%%%%%%%%%%%%%%%%%%%%%%%%%%%%%%%%%%%%%%%%%%%%%%%%%%%%%%%

The classical one-sample KS null distribution assumes that

\[
F_0
\]

is completely specified.

If parameters of \(F_0\) are estimated from the same data, the usual
critical values are generally no longer exact.

Modified procedures or simulation methods may then be required.

%%%%%%%%%%%%%%%%%%%%%%%%%%%%%%%%%%%%%%%%%%%%%%%%%%%%%%%%%%%%
\subsection{Two-Sample Kolmogorov--Smirnov Test}
\label{subsec:two-sample-ks}
%%%%%%%%%%%%%%%%%%%%%%%%%%%%%%%%%%%%%%%%%%%%%%%%%%%%%%%%%%%%

For independent samples with empirical CDFs

\[
F_{n_1}
\]

and

\[
G_{n_2},
\]

the statistic is

\[
\boxed{
D_{n_1,n_2}
=
\sup_x
|
F_{n_1}(x)-G_{n_2}(x)
|.
}
\]

The null hypothesis is

\[
\boxed{
H_0:
F=G.
}
\]

Thus the test compares two entire continuous distributions.

%%%%%%%%%%%%%%%%%%%%%%%%%%%%%%%%%%%%%%%%%%%%%%%%%%%%%%%%%%%%
\subsection{KS Versus Rank-Sum}
\label{subsec:ks-versus-rank-sum}
%%%%%%%%%%%%%%%%%%%%%%%%%%%%%%%%%%%%%%%%%%%%%%%%%%%%%%%%%%%%

The rank-sum test is especially sensitive to systematic ordering or
location-type differences.

The two-sample KS test is sensitive to differences in the entire CDF,
including:

\[
\boxed{
\text{location},
}
\]

\[
\boxed{
\text{spread},
}
\]

or

\[
\boxed{
\text{shape}.
}
\]

Thus the two tests answer different questions.

%%%%%%%%%%%%%%%%%%%%%%%%%%%%%%%%%%%%%%%%%%%%%%%%%%%%%%%%%%%%
\subsection{Cramér--von Mises Statistic Preview}
\label{subsec:cramer-von-mises-preview}
%%%%%%%%%%%%%%%%%%%%%%%%%%%%%%%%%%%%%%%%%%%%%%%%%%%%%%%%%%%%

Rather than using the maximum CDF difference, the Cramér--von Mises
approach integrates squared differences.

Conceptually,

\[
\boxed{
\int
[
F_n(x)-F_0(x)
]^2
\,dF_0(x).
}
\]

Thus discrepancies over the whole distribution contribute to the
statistic.

%%%%%%%%%%%%%%%%%%%%%%%%%%%%%%%%%%%%%%%%%%%%%%%%%%%%%%%%%%%%
\subsection{Anderson--Darling Test Preview}
\label{subsec:anderson-darling-preview}
%%%%%%%%%%%%%%%%%%%%%%%%%%%%%%%%%%%%%%%%%%%%%%%%%%%%%%%%%%%%

The Anderson--Darling approach also compares

\[
F_n
\]

with

\[
F_0,
\]

but places relatively greater weight on tail discrepancies.

This can make it more sensitive than KS to certain tail departures.

%%%%%%%%%%%%%%%%%%%%%%%%%%%%%%%%%%%%%%%%%%%%%%%%%%%%%%%%%%%%
\subsection{Runs Test}
\label{subsec:runs-test}
%%%%%%%%%%%%%%%%%%%%%%%%%%%%%%%%%%%%%%%%%%%%%%%%%%%%%%%%%%%%

A runs test evaluates whether a sequence appears randomly ordered with
respect to two categories.

A run is a maximal consecutive sequence of the same category.

For example,

\[
A\,A\,B\,B\,A\,B\,B
\]

contains four runs:

\[
AA,
\quad
BB,
\quad
A,
\quad
BB.
\]

%%%%%%%%%%%%%%%%%%%%%%%%%%%%%%%%%%%%%%%%%%%%%%%%%%%%%%%%%%%%
\subsection{Expected Number of Runs}
\label{subsec:expected-runs}
%%%%%%%%%%%%%%%%%%%%%%%%%%%%%%%%%%%%%%%%%%%%%%%%%%%%%%%%%%%%

Suppose a sequence contains

\[
n_1
\]

observations of one type and

\[
n_2
\]

of the other.

Under random ordering,

\[
\boxed{
\mathbb{E}[R]
=
1+
\frac{
2n_1n_2
}{
n_1+n_2
}.
}
\]

Very few runs suggest clustering.

Very many runs suggest excessive alternation.

%%%%%%%%%%%%%%%%%%%%%%%%%%%%%%%%%%%%%%%%%%%%%%%%%%%%%%%%%%%%
\subsection{Permutation Tests}
\label{subsec:permutation-tests-nonparametric}
%%%%%%%%%%%%%%%%%%%%%%%%%%%%%%%%%%%%%%%%%%%%%%%%%%%%%%%%%%%%

A permutation test generates a null distribution by rearranging labels
or observations in ways justified by the null hypothesis.

Suppose two groups are compared using statistic

\[
T.
\]

Under an exchangeability null:

\begin{enumerate}
    \item combine the observations;
    \item permute group labels;
    \item recompute \(T\);
    \item repeat over all or many valid permutations;
    \item compare the observed statistic with the permutation
          distribution.
\end{enumerate}

%%%%%%%%%%%%%%%%%%%%%%%%%%%%%%%%%%%%%%%%%%%%%%%%%%%%%%%%%%%%
\subsection{Exact Permutation \(p\)-Value}
\label{subsec:exact-permutation-p-value}
%%%%%%%%%%%%%%%%%%%%%%%%%%%%%%%%%%%%%%%%%%%%%%%%%%%%%%%%%%%%

If all

\[
M
\]

allowable permutations are enumerated, an exact upper-tail permutation
\(p\)-value can be written

\[
\boxed{
p
=
\frac{
\#\{
T^*\geq T_{\text{obs}}
\}
}{
M
}.
}
\]

For a two-sided test, extremeness must be defined according to the chosen
statistic and null hypothesis.

%%%%%%%%%%%%%%%%%%%%%%%%%%%%%%%%%%%%%%%%%%%%%%%%%%%%%%%%%%%%
\subsection{Monte Carlo Permutation Test}
\label{subsec:monte-carlo-permutation}
%%%%%%%%%%%%%%%%%%%%%%%%%%%%%%%%%%%%%%%%%%%%%%%%%%%%%%%%%%%%

When the number of permutations is enormous, randomly sample

\[
B
\]

valid permutations.

A common Monte Carlo estimate is

\[
\boxed{
\widehat{p}
=
\frac{
1+
\#\{
T_b^*
\text{ at least as extreme as }T_{\text{obs}}
\}
}{
B+1
}.
}
\]

The added \(1\) prevents an estimated \(p\)-value of exactly zero and
provides a standard finite-simulation correction.

%%%%%%%%%%%%%%%%%%%%%%%%%%%%%%%%%%%%%%%%%%%%%%%%%%%%%%%%%%%%
\subsection{Permutation Tests and Exchangeability}
\label{subsec:permutation-exchangeability}
%%%%%%%%%%%%%%%%%%%%%%%%%%%%%%%%%%%%%%%%%%%%%%%%%%%%%%%%%%%%

Permutation validity depends on exchangeability under the null.

If observations have unequal distributions or dependence structures that
are not preserved by the permutation, an unrestricted permutation may be
invalid.

Thus:

\[
\boxed{
\text{permutation}
\neq
\text{automatic assumption-free inference}.
}
\]

%%%%%%%%%%%%%%%%%%%%%%%%%%%%%%%%%%%%%%%%%%%%%%%%%%%%%%%%%%%%
\subsection{Randomization Tests Versus Permutation Tests}
\label{subsec:randomization-versus-permutation}
%%%%%%%%%%%%%%%%%%%%%%%%%%%%%%%%%%%%%%%%%%%%%%%%%%%%%%%%%%%%

In randomized experiments, a randomization test uses the actual treatment
assignment mechanism.

Its validity comes directly from the experimental design.

A permutation test in observational data instead relies on a
distributional exchangeability assumption.

The computations may look similar, but their inferential foundations
differ.

%%%%%%%%%%%%%%%%%%%%%%%%%%%%%%%%%%%%%%%%%%%%%%%%%%%%%%%%%%%%
\subsection{Rank Tests as Permutation Tests}
\label{subsec:rank-tests-permutation}
%%%%%%%%%%%%%%%%%%%%%%%%%%%%%%%%%%%%%%%%%%%%%%%%%%%%%%%%%%%%

Many classical rank tests can be understood through permutation
distributions of ranks.

For example, under a two-sample identical-distribution null, ranks are
exchangeable between the two groups.

Thus rank-sum inference can be obtained from the distribution of rank
allocations.

%%%%%%%%%%%%%%%%%%%%%%%%%%%%%%%%%%%%%%%%%%%%%%%%%%%%%%%%%%%%
\subsection{Exact Versus Asymptotic Inference}
\label{subsec:exact-versus-asymptotic-nonparametric}
%%%%%%%%%%%%%%%%%%%%%%%%%%%%%%%%%%%%%%%%%%%%%%%%%%%%%%%%%%%%

Some nonparametric tests have exact finite-sample null distributions.

Examples include:

\[
\boxed{
\text{sign tests},
}
\]

and certain permutation tests.

For larger samples, exact calculations may be computationally expensive.

Then normal or chi-square approximations may be used.

%%%%%%%%%%%%%%%%%%%%%%%%%%%%%%%%%%%%%%%%%%%%%%%%%%%%%%%%%%%%
\subsection{Continuity Corrections}
\label{subsec:continuity-corrections-nonparametric}
%%%%%%%%%%%%%%%%%%%%%%%%%%%%%%%%%%%%%%%%%%%%%%%%%%%%%%%%%%%%

Rank and sign statistics are discrete.

When approximating them using a continuous normal distribution, a
continuity correction may improve finite-sample approximation.

The precise correction depends on the statistic and tail direction.

%%%%%%%%%%%%%%%%%%%%%%%%%%%%%%%%%%%%%%%%%%%%%%%%%%%%%%%%%%%%
\subsection{Tie Corrections}
\label{subsec:tie-corrections-nonparametric}
%%%%%%%%%%%%%%%%%%%%%%%%%%%%%%%%%%%%%%%%%%%%%%%%%%%%%%%%%%%%

Ties alter the variance of rank statistics because ranks are no longer
all distinct.

For example, a group of \(t\) tied observations produces a correction
involving

\[
\boxed{
t^3-t.
}
\]

Software commonly incorporates these corrections automatically.

Exact test distributions may also depend on how ties are handled.

%%%%%%%%%%%%%%%%%%%%%%%%%%%%%%%%%%%%%%%%%%%%%%%%%%%%%%%%%%%%
\subsection{Rank-Based Estimation}
\label{subsec:rank-based-estimation}
%%%%%%%%%%%%%%%%%%%%%%%%%%%%%%%%%%%%%%%%%%%%%%%%%%%%%%%%%%%%

Nonparametric methods are not limited to hypothesis tests.

Rank procedures can also produce point estimates and confidence
intervals.

An important example is the Hodges--Lehmann estimator.

%%%%%%%%%%%%%%%%%%%%%%%%%%%%%%%%%%%%%%%%%%%%%%%%%%%%%%%%%%%%
\subsection{Hodges--Lehmann One-Sample Estimator}
\label{subsec:hodges-lehmann-one-sample}
%%%%%%%%%%%%%%%%%%%%%%%%%%%%%%%%%%%%%%%%%%%%%%%%%%%%%%%%%%%%

For a one-sample symmetric location model, form Walsh averages

\[
\boxed{
W_{ij}
=
\frac{
X_i+X_j
}{2},
\qquad
i\leq j.
}
\]

The Hodges--Lehmann estimator is the median of these Walsh averages.

It is naturally connected to the Wilcoxon signed-rank procedure.

%%%%%%%%%%%%%%%%%%%%%%%%%%%%%%%%%%%%%%%%%%%%%%%%%%%%%%%%%%%%
\subsection{Two-Sample Hodges--Lehmann Estimator}
\label{subsec:hodges-lehmann-two-sample}
%%%%%%%%%%%%%%%%%%%%%%%%%%%%%%%%%%%%%%%%%%%%%%%%%%%%%%%%%%%%

For independent samples

\[
X_1,\ldots,X_m
\]

and

\[
Y_1,\ldots,Y_n,
\]

form all pairwise differences

\[
\boxed{
X_i-Y_j.
}
\]

A Hodges--Lehmann estimator of a location shift is the median of these

\[
mn
\]

pairwise differences.

%%%%%%%%%%%%%%%%%%%%%%%%%%%%%%%%%%%%%%%%%%%%%%%%%%%%%%%%%%%%
\subsection{Rank-Based Confidence Intervals}
\label{subsec:rank-based-confidence-intervals}
%%%%%%%%%%%%%%%%%%%%%%%%%%%%%%%%%%%%%%%%%%%%%%%%%%%%%%%%%%%%

Because signed-rank and rank-sum tests can be inverted over possible
location shifts, they can produce confidence intervals.

The endpoints often correspond to selected ordered Walsh averages or
pairwise differences.

Thus there is a nonparametric analogue of the test--confidence-interval
duality developed earlier.

%%%%%%%%%%%%%%%%%%%%%%%%%%%%%%%%%%%%%%%%%%%%%%%%%%%%%%%%%%%%
\subsection{Robustness}
\label{subsec:nonparametric-robustness}
%%%%%%%%%%%%%%%%%%%%%%%%%%%%%%%%%%%%%%%%%%%%%%%%%%%%%%%%%%%%

Rank-based methods are often resistant to extreme numerical observations.

For example, replacing the largest observation by an even larger value
may leave all ranks unchanged.

Thus the statistic may remain unchanged.

This contrasts with the sample mean and variance, which can change
substantially.

%%%%%%%%%%%%%%%%%%%%%%%%%%%%%%%%%%%%%%%%%%%%%%%%%%%%%%%%%%%%
\subsection{Robust Does Not Mean Universally Better}
\label{subsec:robust-not-universal}
%%%%%%%%%%%%%%%%%%%%%%%%%%%%%%%%%%%%%%%%%%%%%%%%%%%%%%%%%%%%

Nonparametric procedures may sacrifice information by replacing raw
values with ranks or signs.

If a parametric model is accurate, parametric methods can be more
efficient.

Thus the choice involves a tradeoff between:

\[
\boxed{
\text{model efficiency}
}
\]

and

\[
\boxed{
\text{robustness to misspecification}.
}
\]

%%%%%%%%%%%%%%%%%%%%%%%%%%%%%%%%%%%%%%%%%%%%%%%%%%%%%%%%%%%%
\subsection{Asymptotic Relative Efficiency}
\label{subsec:asymptotic-relative-efficiency}
%%%%%%%%%%%%%%%%%%%%%%%%%%%%%%%%%%%%%%%%%%%%%%%%%%%%%%%%%%%%

Asymptotic relative efficiency, abbreviated ARE, compares the sample
sizes needed by two procedures to attain similar asymptotic power.

Schematically,

\[
\boxed{
ARE(T_1,T_2)
=
\lim
\frac{
n_2
}{
n_1
}.
}
\]

An ARE near \(1\) indicates similar large-sample efficiency.

%%%%%%%%%%%%%%%%%%%%%%%%%%%%%%%%%%%%%%%%%%%%%%%%%%%%%%%%%%%%
\subsection{Wilcoxon Efficiency Under Normality}
\label{subsec:wilcoxon-efficiency-normal}
%%%%%%%%%%%%%%%%%%%%%%%%%%%%%%%%%%%%%%%%%%%%%%%%%%%%%%%%%%%%

Under normal location models, the Wilcoxon rank procedures retain high
efficiency relative to classical \(t\)-procedures.

This illustrates that rank methods can be robust without necessarily
sacrificing large amounts of efficiency.

The precise asymptotic efficiency depends on the distributions being
compared.

%%%%%%%%%%%%%%%%%%%%%%%%%%%%%%%%%%%%%%%%%%%%%%%%%%%%%%%%%%%%
\subsection{Breakdown and Resistance}
\label{subsec:breakdown-resistance}
%%%%%%%%%%%%%%%%%%%%%%%%%%%%%%%%%%%%%%%%%%%%%%%%%%%%%%%%%%%%

Robust procedures can also be compared by how much contamination they can
tolerate before becoming arbitrarily unreliable.

The sample median is far more resistant to extreme contamination than the
sample mean.

This connects nonparametric methods with robust statistics.

%%%%%%%%%%%%%%%%%%%%%%%%%%%%%%%%%%%%%%%%%%%%%%%%%%%%%%%%%%%%
\subsection{Quantile Estimation}
\label{subsec:nonparametric-quantile-estimation}
%%%%%%%%%%%%%%%%%%%%%%%%%%%%%%%%%%%%%%%%%%%%%%%%%%%%%%%%%%%%

Population quantiles can be estimated directly from order statistics.

For

\[
0<p<1,
\]

the population \(p\)-quantile satisfies conceptually

\[
F(q_p)\geq p.
\]

A sample quantile is obtained from the ordered sample.

Finite-sample interpolation conventions differ across software.

%%%%%%%%%%%%%%%%%%%%%%%%%%%%%%%%%%%%%%%%%%%%%%%%%%%%%%%%%%%%
\subsection{Sample Median}
\label{subsec:sample-median-nonparametric}
%%%%%%%%%%%%%%%%%%%%%%%%%%%%%%%%%%%%%%%%%%%%%%%%%%%%%%%%%%%%

The median is the

\[
0.5
\]

quantile.

For odd \(n\),

\[
\boxed{
\widetilde{X}
=
X_{((n+1)/2)}.
}
\]

For even \(n\), the conventional descriptive median is often

\[
\boxed{
\widetilde{X}
=
\frac{
X_{(n/2)}
+
X_{(n/2+1)}
}{2}.
}
\]

Other quantile conventions can differ for general percentiles.

%%%%%%%%%%%%%%%%%%%%%%%%%%%%%%%%%%%%%%%%%%%%%%%%%%%%%%%%%%%%
\subsection{Nonparametric Density Estimation}
\label{subsec:nonparametric-density-estimation}
%%%%%%%%%%%%%%%%%%%%%%%%%%%%%%%%%%%%%%%%%%%%%%%%%%%%%%%%%%%%

A probability density can be estimated without specifying a named
parametric family.

One simple estimator is a histogram.

A smoother alternative is kernel density estimation.

%%%%%%%%%%%%%%%%%%%%%%%%%%%%%%%%%%%%%%%%%%%%%%%%%%%%%%%%%%%%
\subsection{Kernel Density Estimator}
\label{subsec:kernel-density-estimator}
%%%%%%%%%%%%%%%%%%%%%%%%%%%%%%%%%%%%%%%%%%%%%%%%%%%%%%%%%%%%

Given observations

\[
X_1,\ldots,X_n,
\]

a kernel density estimator is

\[
\boxed{
\widehat{f}_h(x)
=
\frac1{nh}
\sum_{i=1}^{n}
K
\left(
\frac{
x-X_i
}{h}
\right),
}
\]

where:

\[
K
=
\text{kernel function},
\]

and

\[
h>0
\]

is the bandwidth.

%%%%%%%%%%%%%%%%%%%%%%%%%%%%%%%%%%%%%%%%%%%%%%%%%%%%%%%%%%%%
\subsection{Bandwidth}
\label{subsec:kde-bandwidth}
%%%%%%%%%%%%%%%%%%%%%%%%%%%%%%%%%%%%%%%%%%%%%%%%%%%%%%%%%%%%

The bandwidth controls smoothness.

A small \(h\) produces a highly variable density estimate.

A large \(h\) produces a smoother estimate but can hide important
features.

Thus bandwidth choice reflects a bias--variance tradeoff:

\[
\boxed{
h\text{ small}
\rightarrow
\text{low smoothing, high variance},
}
\]

\[
\boxed{
h\text{ large}
\rightarrow
\text{high smoothing, high bias}.
}
\]

%%%%%%%%%%%%%%%%%%%%%%%%%%%%%%%%%%%%%%%%%%%%%%%%%%%%%%%%%%%%
\subsection{Kernel Functions}
\label{subsec:kernel-functions}
%%%%%%%%%%%%%%%%%%%%%%%%%%%%%%%%%%%%%%%%%%%%%%%%%%%%%%%%%%%%

Common kernels include:

\[
\boxed{
\text{Gaussian},
}
\]

\[
\boxed{
\text{uniform},
}
\]

and

\[
\boxed{
\text{Epanechnikov}.
}
\]

In many applications, bandwidth choice affects the estimate more strongly
than the exact kernel shape.

%%%%%%%%%%%%%%%%%%%%%%%%%%%%%%%%%%%%%%%%%%%%%%%%%%%%%%%%%%%%
\subsection{Nonparametric Regression}
\label{subsec:nonparametric-regression}
%%%%%%%%%%%%%%%%%%%%%%%%%%%%%%%%%%%%%%%%%%%%%%%%%%%%%%%%%%%%

Instead of assuming

\[
\mathbb{E}[Y\mid X=x]
=
\beta_0+\beta_1x,
\]

nonparametric regression estimates a general function

\[
\boxed{
m(x)
=
\mathbb{E}[Y\mid X=x].
}
\]

The function's shape is learned flexibly from data.

%%%%%%%%%%%%%%%%%%%%%%%%%%%%%%%%%%%%%%%%%%%%%%%%%%%%%%%%%%%%
\subsection{Kernel Regression Preview}
\label{subsec:kernel-regression-preview}
%%%%%%%%%%%%%%%%%%%%%%%%%%%%%%%%%%%%%%%%%%%%%%%%%%%%%%%%%%%%

A simple kernel regression estimator is

\[
\boxed{
\widehat{m}(x)
=
\frac{
\sum_{i=1}^{n}
K
\left(
\frac{x-X_i}{h}
\right)
Y_i
}{
\sum_{i=1}^{n}
K
\left(
\frac{x-X_i}{h}
\right)
}.
}
\]

Nearby observations receive greater weight than distant observations.

This is often called the Nadaraya--Watson estimator.

%%%%%%%%%%%%%%%%%%%%%%%%%%%%%%%%%%%%%%%%%%%%%%%%%%%%%%%%%%%%
\subsection{Local Regression Preview}
\label{subsec:local-regression-preview}
%%%%%%%%%%%%%%%%%%%%%%%%%%%%%%%%%%%%%%%%%%%%%%%%%%%%%%%%%%%%

Local polynomial regression fits simple polynomial models near each
target value \(x\).

For example, local linear regression minimizes a weighted sum such as

\[
\boxed{
\sum_{i=1}^{n}
w_i(x)
[
Y_i-a-b(X_i-x)
]^2.
}
\]

The fitted intercept gives the estimated mean response at \(x\).

%%%%%%%%%%%%%%%%%%%%%%%%%%%%%%%%%%%%%%%%%%%%%%%%%%%%%%%%%%%%
\subsection{Smoothing Splines Preview}
\label{subsec:smoothing-splines-preview}
%%%%%%%%%%%%%%%%%%%%%%%%%%%%%%%%%%%%%%%%%%%%%%%%%%%%%%%%%%%%

A smoothing spline balances data fit and smoothness through an objective
of the form

\[
\boxed{
\sum_{i=1}^{n}
[
Y_i-f(X_i)
]^2
+
\lambda
\int
[f''(x)]^2\,dx.
}
\]

The tuning parameter

\[
\lambda
\]

controls smoothness.

This connects nonparametric regression with regularization.

%%%%%%%%%%%%%%%%%%%%%%%%%%%%%%%%%%%%%%%%%%%%%%%%%%%%%%%%%%%%
\subsection{Nonparametric Bootstrap}
\label{subsec:nonparametric-bootstrap}
%%%%%%%%%%%%%%%%%%%%%%%%%%%%%%%%%%%%%%%%%%%%%%%%%%%%%%%%%%%%

The ordinary bootstrap can be interpreted as sampling from the empirical
distribution

\[
F_n.
\]

A bootstrap sample consists of

\[
n
\]

observations drawn with replacement from the original sample.

Thus the bootstrap is naturally nonparametric when no parametric family
is imposed.

%%%%%%%%%%%%%%%%%%%%%%%%%%%%%%%%%%%%%%%%%%%%%%%%%%%%%%%%%%%%
\subsection{Bootstrap Distribution}
\label{subsec:bootstrap-distribution-nonparametric}
%%%%%%%%%%%%%%%%%%%%%%%%%%%%%%%%%%%%%%%%%%%%%%%%%%%%%%%%%%%%

For statistic

\[
T_n
=
T(X_1,\ldots,X_n),
\]

generate bootstrap samples and compute

\[
T_1^*,
\ldots,
T_B^*.
\]

Their empirical distribution approximates the sampling distribution of

\[
T_n
\]

under suitable conditions.

%%%%%%%%%%%%%%%%%%%%%%%%%%%%%%%%%%%%%%%%%%%%%%%%%%%%%%%%%%%%
\subsection{Permutation Versus Bootstrap}
\label{subsec:permutation-versus-bootstrap}
%%%%%%%%%%%%%%%%%%%%%%%%%%%%%%%%%%%%%%%%%%%%%%%%%%%%%%%%%%%%

Permutation methods primarily approximate a null distribution by
rearranging observations or labels.

Bootstrap methods approximate sampling variability by resampling with
replacement.

Thus:

\[
\boxed{
\text{permutation}
\rightarrow
\text{testing under exchangeability},
}
\]

while

\[
\boxed{
\text{bootstrap}
\rightarrow
\text{sampling-distribution approximation}.
}
\]

%%%%%%%%%%%%%%%%%%%%%%%%%%%%%%%%%%%%%%%%%%%%%%%%%%%%%%%%%%%%
\subsection{Choosing a Nonparametric Procedure}
\label{subsec:choosing-nonparametric-procedure}
%%%%%%%%%%%%%%%%%%%%%%%%%%%%%%%%%%%%%%%%%%%%%%%%%%%%%%%%%%%%

For one sample or paired data and a median-type question:

\[
\boxed{
\text{sign test}.
}
\]

For paired symmetric location differences:

\[
\boxed{
\text{Wilcoxon signed-rank test}.
}
\]

For two independent groups:

\[
\boxed{
\text{Mann--Whitney--Wilcoxon rank-sum test}.
}
\]

For several independent groups:

\[
\boxed{
\text{Kruskal--Wallis test}.
}
\]

For blocked or repeated treatments:

\[
\boxed{
\text{Friedman test}.
}
\]

For monotonic association:

\[
\boxed{
\text{Spearman or Kendall correlation}.
}
\]

For comparing full distributions:

\[
\boxed{
\text{Kolmogorov--Smirnov procedures}.
}
\]

%%%%%%%%%%%%%%%%%%%%%%%%%%%%%%%%%%%%%%%%%%%%%%%%%%%%%%%%%%%%
\subsection{Parametric Versus Nonparametric Choice}
\label{subsec:parametric-versus-nonparametric-choice}
%%%%%%%%%%%%%%%%%%%%%%%%%%%%%%%%%%%%%%%%%%%%%%%%%%%%%%%%%%%%

The decision should not be based only on whether a normality test rejects.

Consider:

\[
\boxed{
\text{scientific parameter of interest},
}
\]

\[
\boxed{
\text{measurement scale},
}
\]

\[
\boxed{
\text{sample size},
}
\]

\[
\boxed{
\text{outliers},
}
\]

\[
\boxed{
\text{distributional shape},
}
\]

and

\[
\boxed{
\text{study design}.
}
\]

A nonparametric test may answer a different scientific question from its
parametric counterpart.

%%%%%%%%%%%%%%%%%%%%%%%%%%%%%%%%%%%%%%%%%%%%%%%%%%%%%%%%%%%%
\subsection{Nonparametric Does Not Mean Median}
\label{subsec:nonparametric-not-median}
%%%%%%%%%%%%%%%%%%%%%%%%%%%%%%%%%%%%%%%%%%%%%%%%%%%%%%%%%%%%

A common mistake is to interpret every rank test as testing medians.

Different procedures target different distributional features.

For example:

\[
\boxed{
\text{sign test}
}
\]

can directly address a median under suitable continuity conditions,

while

\[
\boxed{
\text{Mann--Whitney}
}
\]

more generally concerns relative distributional ordering.

Interpretation must follow the actual null hypothesis.

%%%%%%%%%%%%%%%%%%%%%%%%%%%%%%%%%%%%%%%%%%%%%%%%%%%%%%%%%%%%
\subsection{Common Errors}
\label{subsec:nonparametric-common-errors}
%%%%%%%%%%%%%%%%%%%%%%%%%%%%%%%%%%%%%%%%%%%%%%%%%%%%%%%%%%%%

\subsubsection{Calling Nonparametric Methods Assumption-Free}

Nonparametric procedures still rely on assumptions such as independence,
exchangeability, symmetry, or continuity.

%%%%%%%%%%%%%%%%%%%%%%%%%%%%%%%%%%%%%%%%%%%%%%%%%%%%%%%%%%%%
\subsubsection{Using an Independent-Sample Rank Test for Paired Data}

Paired observations require a paired procedure such as the sign test or
signed-rank test.

Ignoring pairing loses information and gives the wrong null distribution.

%%%%%%%%%%%%%%%%%%%%%%%%%%%%%%%%%%%%%%%%%%%%%%%%%%%%%%%%%%%%
\subsubsection{Using a Paired Test for Independent Groups}

Independent groups do not provide meaningful within-pair differences.

Use a two-independent-sample procedure instead.

%%%%%%%%%%%%%%%%%%%%%%%%%%%%%%%%%%%%%%%%%%%%%%%%%%%%%%%%%%%%
\subsubsection{Interpreting Mann--Whitney Automatically as a Median Test}

Different distributional shapes can produce a significant rank-sum test
even when medians are similar.

Median interpretation requires additional location-shift assumptions.

%%%%%%%%%%%%%%%%%%%%%%%%%%%%%%%%%%%%%%%%%%%%%%%%%%%%%%%%%%%%
\subsubsection{Ignoring Symmetry for Signed-Rank Location Inference}

The signed-rank procedure uses both signs and absolute ranks.

Its classical location interpretation assumes symmetry of the difference
distribution.

%%%%%%%%%%%%%%%%%%%%%%%%%%%%%%%%%%%%%%%%%%%%%%%%%%%%%%%%%%%%
\subsubsection{Ignoring Ties}

Ties alter rank assignments and may change null variances.

Appropriate tie corrections or exact methods should be used.

%%%%%%%%%%%%%%%%%%%%%%%%%%%%%%%%%%%%%%%%%%%%%%%%%%%%%%%%%%%%
\subsubsection{Applying Standard KS Critical Values After Estimating Parameters}

If a reference distribution's parameters are estimated using the same
data, the standard one-sample KS null distribution generally changes.

%%%%%%%%%%%%%%%%%%%%%%%%%%%%%%%%%%%%%%%%%%%%%%%%%%%%%%%%%%%%
\subsubsection{Assuming KS Tests Only Location}

KS tests compare entire empirical CDFs.

A significant result may arise from differences in location, scale, or
shape.

%%%%%%%%%%%%%%%%%%%%%%%%%%%%%%%%%%%%%%%%%%%%%%%%%%%%%%%%%%%%
\subsubsection{Permuting Labels Without Exchangeability}

A permutation distribution is valid only when the proposed rearrangements
are justified under the null or experimental design.

%%%%%%%%%%%%%%%%%%%%%%%%%%%%%%%%%%%%%%%%%%%%%%%%%%%%%%%%%%%%
\subsubsection{Using Raw Observations Instead of Within-Block Ranks in Friedman Testing}

The Friedman procedure ranks treatments separately within each block.

It does not rank all observations globally.

%%%%%%%%%%%%%%%%%%%%%%%%%%%%%%%%%%%%%%%%%%%%%%%%%%%%%%%%%%%%
\subsubsection{Ignoring Multiple Comparisons After a Significant Omnibus Test}

A significant Kruskal--Wallis or Friedman statistic does not identify the
specific groups or treatments that differ.

Follow-up comparisons require multiplicity control.

%%%%%%%%%%%%%%%%%%%%%%%%%%%%%%%%%%%%%%%%%%%%%%%%%%%%%%%%%%%%
\subsubsection{Assuming Rank Tests Are Immune to Outliers in Every Sense}

Ranks reduce sensitivity to extreme magnitudes, but unusual observations
can still affect ranks, group ordering, and interpretation.

%%%%%%%%%%%%%%%%%%%%%%%%%%%%%%%%%%%%%%%%%%%%%%%%%%%%%%%%%%%%
\subsubsection{Choosing a Test Solely Because a Normality Test Is Significant}

Formal normality tests can reject negligible departures in large samples
and have little power in small samples.

Graphical diagnostics, robustness, sample size, and the scientific target
should also be considered.

%%%%%%%%%%%%%%%%%%%%%%%%%%%%%%%%%%%%%%%%%%%%%%%%%%%%%%%%%%%%
\subsection{Applications}
\label{subsec:nonparametric-applications}
%%%%%%%%%%%%%%%%%%%%%%%%%%%%%%%%%%%%%%%%%%%%%%%%%%%%%%%%%%%%

\begin{applicationbox}

\textbf{Medicine.}
Rank and sign methods analyze skewed biomarkers, ordinal outcomes,
small clinical samples, and paired measurements.

\medskip

\textbf{Environmental Science.}
Environmental concentrations often contain skewness, detection limits,
and extreme observations that motivate robust and rank-based procedures.

\medskip

\textbf{Education Research.}
Ordinal survey responses, small program samples, and paired student
outcomes may be analyzed using rank procedures.

\medskip

\textbf{Psychology.}
Rank correlations and repeated-measures rank tests are useful for ordinal
and nonnormally distributed behavioral data.

\medskip

\textbf{Engineering.}
Nonparametric procedures provide robust comparisons when lifetime or
failure measurements do not follow simple parametric distributions.

\medskip

\textbf{Business.}
Customer ratings, rankings, spending distributions, and A/B experiments
can be studied with rank and permutation procedures.

\medskip

\textbf{Computational Statistics.}
Permutation methods, bootstrap procedures, and empirical distributions
are central computational tools.

\medskip

\textbf{Machine Learning.}
Kernel methods, nonparametric regression, density estimation, and
rank-based performance measures extend statistical modeling beyond fixed
parametric forms.

\end{applicationbox}

%%%%%%%%%%%%%%%%%%%%%%%%%%%%%%%%%%%%%%%%%%%%%%%%%%%%%%%%%%%%
\subsection{Exercises}
\label{subsec:nonparametric-exercises}
%%%%%%%%%%%%%%%%%%%%%%%%%%%%%%%%%%%%%%%%%%%%%%%%%%%%%%%%%%%%

\begin{exercise}[1]
Explain the difference between parametric and nonparametric statistics.
\end{exercise}

\begin{exercise}[1]
Explain why nonparametric does not mean assumption-free.
\end{exercise}

\begin{exercise}[1]
Define a rank.
\end{exercise}

\begin{exercise}[1]
Define the empirical distribution function.
\end{exercise}

\begin{exercise}[1]
State the null distribution of the one-sample sign-test statistic.
\end{exercise}

\begin{exercise}[1]
Describe the Wilcoxon signed-rank procedure.
\end{exercise}

\begin{exercise}[1]
Describe the Wilcoxon rank-sum procedure.
\end{exercise}

\begin{exercise}[1]
Define the Mann--Whitney \(U\)-statistic.
\end{exercise}

\begin{exercise}[1]
State the Kruskal--Wallis statistic.
\end{exercise}

\begin{exercise}[1]
State the Friedman statistic.
\end{exercise}

\begin{exercise}[1]
Define Spearman rank correlation.
\end{exercise}

\begin{exercise}[1]
Define Kendall's tau.
\end{exercise}

\begin{exercise}[1]
Define the one-sample Kolmogorov--Smirnov statistic.
\end{exercise}

\begin{exercise}[1]
Define the two-sample Kolmogorov--Smirnov statistic.
\end{exercise}

\begin{exercise}[2]
Rank the observations

\[
4,\ 10,\ 7,\ 2,\ 8.
\]
\end{exercise}

\begin{exercise}[2]
Assign average ranks to

\[
2,\ 5,\ 5,\ 5,\ 9.
\]
\end{exercise}

\begin{exercise}[2]
Suppose \(15\) observations are compared with a hypothesized median and
\(11\) lie above it. Write the exact right-tailed sign-test \(p\)-value.
\end{exercise}

\begin{exercise}[2]
Suppose paired differences are

\[
-4,\ 2,\ 1,\ -3.
\]

Compute \(W^+\) and \(W^-\).
\end{exercise}

\begin{exercise}[2]
For \(n=10\), compute

\[
\mathbb{E}[W^+].
\]
\end{exercise}

\begin{exercise}[2]
For \(n=10\), compute the no-tie variance of \(W^+\).
\end{exercise}

\begin{exercise}[2]
Suppose group \(1\) has \(n_1=5\), group \(2\) has \(n_2=7\), and the
rank sum for group \(1\) is \(R_1=25\). Compute \(U_1\).
\end{exercise}

\begin{exercise}[2]
For the previous problem, compute \(U_2\).
\end{exercise}

\begin{exercise}[2]
Suppose \(N=20\), \(n_1=8\). Compute the null expected rank sum for group
\(1\).
\end{exercise}

\begin{exercise}[2]
Three groups have rank sums

\[
20,\ 30,\ 40
\]

and sample sizes

\[
5,\ 5,\ 5.
\]

Write the Kruskal--Wallis statistic.
\end{exercise}

\begin{exercise}[2]
A Friedman design has \(b=8\) blocks and \(k=4\) treatments. What are the
approximate chi-square degrees of freedom?
\end{exercise}

\begin{exercise}[2]
Compute Spearman's correlation when

\[
\sum d_i^2=10
\]

and

\[
n=5.
\]
\end{exercise}

\begin{exercise}[2]
For \(n=6\), determine the total number of pairs used in Kendall's tau.
\end{exercise}

\begin{exercise}[2]
If \(C=12\) and \(D=3\) among \(15\) untied pairs, compute Kendall's tau.
\end{exercise}

\begin{exercise}[3]
Derive the null binomial distribution of the sign-test statistic.
\end{exercise}

\begin{exercise}[3]
Explain how order statistics can produce an exact confidence interval for
a population median.
\end{exercise}

\begin{exercise}[3]
Derive

\[
W^++W^-
=
\frac{
n(n+1)
}{2}.
\]
\end{exercise}

\begin{exercise}[3]
Derive the null expectation of \(W^+\).
\end{exercise}

\begin{exercise}[3]
Derive the no-tie variance of \(W^+\).
\end{exercise}

\begin{exercise}[3]
Explain why symmetry is important for the classical signed-rank location
interpretation.
\end{exercise}

\begin{exercise}[3]
Derive the relation between the Wilcoxon rank sum and Mann--Whitney
\(U\)-statistic.
\end{exercise}

\begin{exercise}[3]
Show that

\[
U_1+U_2=n_1n_2.
\]
\end{exercise}

\begin{exercise}[3]
Explain the probabilistic interpretation of

\[
\frac{U}{n_1n_2}.
\]
\end{exercise}

\begin{exercise}[3]
Explain why Mann--Whitney does not automatically test equality of
medians.
\end{exercise}

\begin{exercise}[3]
Derive the no-tie expected rank sum under the two-sample null.
\end{exercise}

\begin{exercise}[3]
Derive the Kruskal--Wallis statistic from deviations of group mean ranks
from the overall mean rank.
\end{exercise}

\begin{exercise}[3]
Explain why Kruskal--Wallis is the rank analogue of one-way ANOVA.
\end{exercise}

\begin{exercise}[3]
Explain why Friedman ranks within blocks.
\end{exercise}

\begin{exercise}[3]
Show that the Spearman shortcut formula follows from Pearson correlation
of untied ranks.
\end{exercise}

\begin{exercise}[3]
Explain the concordance interpretation of Kendall's tau.
\end{exercise}

\begin{exercise}[3]
Construct two datasets having perfect Spearman correlation but a
nonlinear relationship.
\end{exercise}

\begin{exercise}[3]
Explain why the two-sample KS test can detect scale differences even when
the population medians are equal.
\end{exercise}

\begin{exercise}[3]
Explain why estimating parameters from the sample changes the classical
one-sample KS null distribution.
\end{exercise}

\begin{exercise}[3]
Compare KS, Cramér--von Mises, and Anderson--Darling conceptually.
\end{exercise}

\begin{exercise}[3]
Derive the expected number of runs for a two-category random sequence.
\end{exercise}

\begin{exercise}[3]
Explain why permutation testing requires exchangeability.
\end{exercise}

\begin{exercise}[3]
Distinguish permutation tests from bootstrap procedures.
\end{exercise}

\begin{exercise}[3]
Explain why a randomization test can be exact under a randomized
experimental design.
\end{exercise}

\begin{exercise}[3]
Explain why ties require variance corrections in rank tests.
\end{exercise}

\begin{exercise}[3]
Construct the Walsh averages for the sample

\[
1,\ 3,\ 6.
\]
\end{exercise}

\begin{exercise}[3]
Find the two-sample Hodges--Lehmann estimator for

\[
X=(2,4),
\qquad
Y=(1,3).
\]
\end{exercise}

\begin{exercise}[3]
Explain the bandwidth bias--variance tradeoff in kernel density
estimation.
\end{exercise}

\begin{exercise}[3]
Explain how kernel regression differs from ordinary linear regression.
\end{exercise}

\begin{exercise}[4]
Prove the Glivenko--Cantelli theorem.
\end{exercise}

\begin{exercise}[4]
Study the Dvoretzky--Kiefer--Wolfowitz inequality.
\end{exercise}

\begin{exercise}[4]
Derive the exact distribution of the sign-test statistic.
\end{exercise}

\begin{exercise}[4]
Derive the exact null distribution of the Wilcoxon signed-rank statistic.
\end{exercise}

\begin{exercise}[4]
Derive the exact permutation distribution of the rank-sum statistic.
\end{exercise}

\begin{exercise}[4]
Study Pitman asymptotic relative efficiency of Wilcoxon and \(t\)-tests.
\end{exercise}

\begin{exercise}[4]
Investigate the asymptotic distribution of the Kruskal--Wallis statistic.
\end{exercise}

\begin{exercise}[4]
Investigate exact Friedman tests.
\end{exercise}

\begin{exercise}[4]
Derive asymptotic normality of Spearman rank correlation.
\end{exercise}

\begin{exercise}[4]
Study Kendall's tau as a \(U\)-statistic.
\end{exercise}

\begin{exercise}[4]
Derive the asymptotic null distribution of the one-sample KS statistic.
\end{exercise}

\begin{exercise}[4]
Study the Kolmogorov distribution.
\end{exercise}

\begin{exercise}[4]
Investigate two-sample empirical-process convergence.
\end{exercise}

\begin{exercise}[4]
Study Anderson--Darling goodness-of-fit statistics.
\end{exercise}

\begin{exercise}[4]
Investigate permutation tests under heteroskedasticity.
\end{exercise}

\begin{exercise}[4]
Study studentized permutation statistics.
\end{exercise}

\begin{exercise}[4]
Study rank-based multiple comparisons after Kruskal--Wallis.
\end{exercise}

\begin{exercise}[4]
Investigate rank-based estimators and confidence intervals.
\end{exercise}

\begin{exercise}[4]
Derive bias and variance properties of kernel density estimators.
\end{exercise}

\begin{exercise}[4]
Study bandwidth selection by cross-validation.
\end{exercise}

\begin{exercise}[4]
Derive properties of local polynomial regression.
\end{exercise}

\begin{exercise}[4]
Study smoothing splines through penalized least squares.
\end{exercise}

%%%%%%%%%%%%%%%%%%%%%%%%%%%%%%%%%%%%%%%%%%%%%%%%%%%%%%%%%%%%
\subsection{Conceptual Summary}
\label{subsec:nonparametric-summary}
%%%%%%%%%%%%%%%%%%%%%%%%%%%%%%%%%%%%%%%%%%%%%%%%%%%%%%%%%%%%

Nonparametric statistics reduces reliance on fully specified parametric
distribution families.

The empirical distribution function is

\[
\boxed{
F_n(x)
=
\frac1n
\sum_{i=1}^{n}
\mathbf{1}_{\{X_i\leq x\}}.
}
\]

For a one-sample median test, the sign statistic satisfies

\[
\boxed{
S
\sim
\operatorname{Binomial}
\left(
n,\frac12
\right)
}
\]

under the continuous median null.

The Wilcoxon signed-rank procedure uses signed ranks

\[
\boxed{
W^+
=
\sum_{D_i>0}R_i.
}
\]

For two independent groups, the Mann--Whitney statistic can be written

\[
\boxed{
U_1
=
R_1
-
\frac{
n_1(n_1+1)
}{2}.
}
\]

The Kruskal--Wallis statistic is

\[
\boxed{
H
=
\frac{
12
}{
N(N+1)
}
\sum_i
\frac{
R_i^2
}{
n_i
}
-
3(N+1).
}
\]

For blocked rank data, Friedman uses

\[
\boxed{
Q
=
\frac{
12
}{
bk(k+1)
}
\sum_jR_j^2
-
3b(k+1).
}
\]

Spearman correlation measures correlation between ranks, while Kendall's
tau measures pairwise concordance.

The Kolmogorov--Smirnov statistic compares entire distribution functions:

\[
\boxed{
D_n
=
\sup_x
|F_n(x)-F_0(x)|.
}
\]

Permutation tests generate null distributions by valid rearrangements
under exchangeability or a randomized design.

Nonparametric statistics therefore combines

\[
\boxed{
\text{order statistics}
+
\text{ranks}
+
\text{empirical distributions}
+
\text{combinatorics}
+
\text{resampling}.
}
\]

%%%%%%%%%%%%%%%%%%%%%%%%%%%%%%%%%%%%%%%%%%%%%%%%%%%%%%%%%%%%
\subsection{Section Connections}
\label{subsec:nonparametric-connections}
%%%%%%%%%%%%%%%%%%%%%%%%%%%%%%%%%%%%%%%%%%%%%%%%%%%%%%%%%%%%

\chapterconnection{
Nonparametric Statistics complements the parametric inference developed
in Sections~\ref{sec:statistical-estimation},
\ref{sec:confidence-intervals}, and
\ref{sec:hypothesis-testing}. Rank-sum and Kruskal--Wallis procedures
provide alternatives to classical mean comparisons, while signed-rank
and Friedman methods extend the same ideas to paired and blocked data.
Permutation methods connect directly with the randomization principles
of Section~\ref{sec:experimental-design}. The empirical distribution,
kernel density estimation, and nonparametric regression also lead toward
the more flexible high-dimensional methods studied later in Statistical
Learning and Computational Statistics. The next section, Multivariate
Statistics, studies datasets in which several quantitative variables are
analyzed jointly rather than one response at a time.
}

%%%%%%%%%%%%%%%%%%%%%%%%%%%%%%%%%%%%%%%%%%%%%%%%%%%%%%%%%%%%
% SECTION SOURCE TRACKING
%%%%%%%%%%%%%%%%%%%%%%%%%%%%%%%%%%%%%%%%%%%%%%%%%%%%%%%%%%%%

%------------------------------------------------------------
% 8.11 Nonparametric Statistics
%------------------------------------------------------------
%
% Source basis:
%   - Standard nonparametric statistical theory.
%   - Standard rank-based inference.
%   - Standard permutation and empirical-distribution methods.
%
% Used for:
%   - parametric versus nonparametric inference;
%   - ranks and ties;
%   - signs;
%   - empirical distribution functions;
%   - Glivenko--Cantelli preview;
%   - sign tests;
%   - distribution-free median intervals;
%   - Wilcoxon signed-rank tests;
%   - sign versus signed-rank procedures;
%   - rank-sum tests;
%   - Mann--Whitney U statistics;
%   - probabilistic interpretation of U;
%   - Kruskal--Wallis tests;
%   - Friedman tests;
%   - Spearman rank correlation;
%   - Kendall rank correlation;
%   - Kolmogorov--Smirnov tests;
%   - Cramér--von Mises preview;
%   - Anderson--Darling preview;
%   - runs tests;
%   - permutation tests;
%   - randomization tests;
%   - exact and asymptotic inference;
%   - continuity corrections;
%   - tie corrections;
%   - Hodges--Lehmann estimators;
%   - rank-based confidence intervals;
%   - robustness;
%   - asymptotic relative efficiency;
%   - quantile estimation;
%   - kernel density estimation;
%   - nonparametric regression;
%   - local regression;
%   - smoothing splines;
%   - nonparametric bootstrap;
%   - procedure selection;
%   - applications and exercises.
%
% Notes:
%   - Multivariate Statistics is developed next in Section 8.12.
%   - Mathematical Statistics later develops empirical-process,
%     U-statistic, and asymptotic theory more deeply.
%   - Statistical Learning later develops kernel, smoothing, and flexible
%     predictive methods.
%   - Computational Statistics develops bootstrap, permutation, and
%     resampling procedures more fully.
%
%%%%%%%%%%%%%%%%%%%%%%%%%%%%%%%%%%%%%%%%%%%%%%%%%%%%%%%%%%%%